\subsection{Evidence verification and ordering}
\label{sec:e6}\label{sec:inflation}\label{sec:eval-verify}
Verification and ranking are distinct operations. One archived industrial
hybrid evaluation recorded 191 candidate scores over 16 runs: 17 in-pool golds
and 174 non-golds. All recorded scores exceeded the configured 0.5 threshold;
one additional candidate had no recorded score. The gold and non-gold medians
were 0.88 and 0.86. The candidate-level Mann--Whitney result ($p=0.26$) is
exploratory because candidates from a question are clustered; it does not
establish equality of the score distributions.

A separate archived same-pool replay compared ranking orders for 16 questions
in the historical hybrid wrapper. Hit@5 was 15/16 when the recorded survivors
retained fusion order and 7/16 when ordered by verifier score
(Figure~\ref{fig:inflation}). The score distribution and the ordering replay
come from distinct archived artifacts, not necessarily the same individual
requests. The observation supports retaining the retrieval order in this
historical setting; it does not establish the benefit of threshold filtering
or a general result for all corpora. NLI and span-level replays returned
5/16 and 3/16, respectively, without isolating the cause of loss. The arXiv
corpus did not show the same collapse, but corpus, domain and language changed
together.

\begin{figure*}[t]
\centering
\includegraphics[width=\textwidth]{fig_exp_verifier}
\caption{Historical verifier diagnostics from separate industrial-corpus
artifacts. (a)~The 191 recorded candidate scores from 16 hybrid runs include
17 gold and 174 non-gold candidates; one additional candidate has no recorded
score. Horizontal bars mark medians; the dashed line marks the 0.5 threshold.
(b)~A separate 16-question same-pool replay compares Hit@5 when recorded
survivors retain fusion order or are reordered by verifier score. Filled and
open symbols denote hits and misses, respectively. The panels need not
represent identical requests; the replay diagnoses ordering, not an isolated
filtering benefit.}
\label{fig:inflation}
\end{figure*}

\subsection{Answer behaviour and operational boundaries}
\label{sec:e5}\label{sec:agentic}\label{sec:abstention}
\label{sec:composition}\label{sec:eval-agent}\label{sec:sysarms}\label{sec:e7}
Table~\ref{tab:e5-answers} summarises answer-level proxies separately for each
run. The original ten-question grid compares the platform agent (P), direct
file tools (F) and a vector-only arm (C); the latter used an 800-character
chunk index. Gold IDs appeared in the answer or trace for all ten cells of
each arm. Strict keyword coverage was 9/10 in all arms, while answers with
recognised citation markers numbered 9/10, 7/10 and 2/10 for P, F and C.
Those numbers count \emph{answers with markers}, not the total number of
markers; question-level results appear in Supplementary Table S16. A later ten-question grid instructed C to search the five selected
libraries directly and merge results, rather than use the chunk index. In that
grid strict coverage was 7/10, 7/10 and 8/10 and marker-bearing answers were
9/10 for each arm. Because the C substrate and instructions changed, the two
grids are not a controlled replication of that arm or proof of an interface
effect. The detailed source-presence audit in Supplementary Tables S7--S8 is
not an entailment test.

\begin{table}[t]
\centering
\small
\setlength{\tabcolsep}{3pt}
\caption{Answer and operational observations from distinct runs. P, F and C
are platform-agent, file-tool and vector-only arms. Answer counts have ten
questions per arm; citation-like markers do not establish source support.}
\label{tab:e5-answers}
\begin{tabular}{@{}p{0.26\columnwidth}p{0.68\columnwidth}@{}}
\toprule
Run & Recorded result and boundary \\
\midrule
Original P/F/C & Strict keywords 9/9/9; marker-bearing answers 9/7/2, each out of 10. Gold ID in answer or trace: 10/10 per arm. \\
Later P/F/C & Strict keywords 7/7/8; marker-bearing answers 9/9/9, each out of 10. Vector arm's substrate and instructions changed. \\
Nov. 15 answers & Gold ID in answer 10/10; all keywords 9/10; recognised marker 10/10. Different scoring rules; one run. \\
Harness case & Eight responses on four questions; two harnesses, non-blinded case review. \\
Concurrency & 111/270 timeouts at 30 clients; one instance, no complete sweep after repair. \\
Index coverage & Eight affected collections in one inspection; scoped and global probes could fail. \\
\bottomrule
\end{tabular}
\end{table}

In the separate 15 November platform run, each of ten answers contained its
gold ID, nine contained all configured keywords, and ten contained a
recognised marker. The full-keyword and marker rules differ from those of the
historical grid, and markers may include the gold identifier itself. These
observations are neither independent semantic-quality measures nor a matched
comparison with the historical arms. The ten API-recorded latencies averaged
128.94 s per question and the recorded total cost was \$3.105. The separate
scripted retrieval diagnostic can fall back to a gold ID and did not retain
full judge outputs; it cannot establish that all golds survived verification.

A separate DeepSeek-R1 demonstration used two agent harnesses on four
questions. A non-blinded case review assessed all eight responses against the
source on answer content, provenance, quotations and fabricated facts. Both
responses to an electricity-cost question reported that the requested figure
was absent while distinguishing a GPU-rental-based training-cost estimate.
This is an observed case of abstention, not a general refusal guarantee.
The case details appear in Supplementary Table S15.

A 27-cell stability campaign had no recorded answer errors across three tasks,
three lanes and three rounds. In a single-instance concurrency sweep of
90/180/270 requests at 5/15/30 clients, median latency rose from 5.2 to
1577 ms and 111/270 requests timed out at 30 clients. After a cache repair,
selected 30-client read-ladder median latency was 65.8 ms; the full sweep was
not repeated. Raw concurrent-write counts and crash-recovery evidence were
unavailable. In a separate inspection, unreadable index segments caused
collections to be skipped without a degraded response marker; logs identified
eight affected collections. A scoped query found a newly ingested document
while the global query missed it, and forced re-index restored the global hit.
Another affected library also had an empty scoped query, so the fault is not
confined to global search.

A reconstructed fixed-candidate gate comparison used ten questions. The three
arms packed the gold first for eight answerable questions, and filtered and
unfiltered document sets matched in nine of ten. Grades varied even for
identical packs; the two non-gold questions were abstained on in all arms.
These small observations do not establish an answer-quality benefit from the
gate. Likewise, a deliberately wrong pinned library yielded zero Hit@5 in a
scope diagnostic, but unscoped search covered additional live libraries and
was not matched to the evaluation corpus. Together, the results establish
selected functional behaviour and operational limitations under documented
conditions, while leaving general answer correctness, deployment reliability
and a causal advantage of tiered navigation open.
